% 388_CMB_Grundform_De_ch.tex
% Johann Pascher, 2. Okt. 2026
% Die CMB-Temperatur in der Grundform: eV-Relation, Verhältnis zu m_e, H0-Formen

\providecommand{\BM}{\textbf{[B]}}
\providecommand{\KM}{\textbf{[K]}}
\providecommand{\SM}{\textbf{[S]}}
\providecommand{\XM}{\textbf{[X]}}
\providecommand{\QM}{\textbf{[Q]}}

\chapter*{Die CMB-Temperatur in der Grundform}
\addcontentsline{toc}{chapter}{Die CMB-Temperatur in der Grundform}

\begin{abstract}
Die Casimir-CMB-Verbindung ist eine Identität, weil die Länge $L_\xi$ allein über die
CMB-Temperatur festgelegt ist (R136). Dieses Dokument stellt alle Wege zusammen, die
$T_{\mathrm{CMB}}$ ohne diese Zirkularität aus $\xi$ gewinnen sollen, und rechnet sie in der
Grundform nach: natürliche Einheiten mit $\alpha=1$, ein einziger Anker $m_e$. Die eV-Relation
$T_{\mathrm{CMB}}=\frac{16}{9}\xi$~eV mit Korrektur zweiter Ordnung trifft auf $2\times10^{-6}$,
gilt aber nur in der Einheit eV, deren Größe 1881 festgelegt und 2019 exakt gesetzt wurde. In
der Grundform hängt $T_{\mathrm{CMB}}/m_e$ an einer einzigen Zahl. Der Kandidat
$T_{\mathrm{CMB}}/m_e=(8\pi)^{1/4}\,\xi^{5/2}$ trifft den FIRAS-Wert auf $0{,}12\sigma$, ergibt mit
der $H_0$-Form des Korpus $T_{\mathrm{CMB}}^4=16\,H_0\,m_e^3$ und macht $L_\xi$ unabhängig von der
CMB. Die $H_0$-Form mit Exponent~10 passt dazu, die ältere mit $41/4$ liegt $11\sigma$ daneben.
Hergeleitet sind weder der Vorfaktor noch der Exponent; was offen bleibt, steht am Ende.
\end{abstract}

% ============================================================
\section*{Statusmarkierungen}
% ============================================================
\begin{center}
\begin{tabular}{@{}lp{11cm}@{}}
\textbf{[B]} & Algebraisch bewiesen. \\[2pt]
\textbf{[K]} & Numerisch bestätigt (Prüfskript, Vergleich mit Messwerten). \\[2pt]
\textbf{[S]} & Setzung, offen oder spekulativ. \\[2pt]
\textbf{[X]} & Widerlegt oder nicht tragfähig. \\[2pt]
\textbf{[Q]} & Konvention. \\
\end{tabular}
\end{center}
Messwerte: $T_{\mathrm{CMB}}=2{,}72548\pm0{,}00057$~K (FIRAS, Fixsen 2009; relative Unsicherheit
$0{,}021\,\%$), CODATA 2022. Prüfskript: \texttt{pruef\_388\_cmb\_grundform.py}.

% ============================================================
\section{Ausgangslage}
% ============================================================

Die Länge $L_\xi=(15\xi/\pi^2)^{1/4}\,\hbar c/(k_BT_{\mathrm{CMB}})=100{,}24\,\mu$m ist so
gewählt, dass die Photonendichte der CMB über ihr den Wert $\xi/L_\xi^4$ hat. Der Vergleich mit
der Casimir-Energiedichte bei derselben Länge ist deshalb eine Identität, keine Messung
(Dok.~279, R136). Auch ein $H_0$-Anker ändert daran nichts, denn $L_\xi$ hängt nur an
$T_{\mathrm{CMB}}$.

Das $\Lambda$CDM-Modell leitet $T_{\mathrm{CMB}}$ ebenfalls nicht her. Es übernimmt den
FIRAS-Wert als feste Eingabe, und der Planck-Wert $H_0=67{,}4$~km/s/Mpc ist eine abgeleitete Größe,
die diesen Wert voraussetzt. Die Frage ist also nicht, ob die FFGFT schlechter dasteht als
$\Lambda$CDM, sondern ob sie mehr leisten kann: $T_{\mathrm{CMB}}$ als Verhältnis aus $\xi$, mit
demselben Anker wie die Massen.

% ============================================================
\section{Die eV-Relation}
% ============================================================

Im Korpus steht seit 2025 die Relation $T_{\mathrm{CMB}}=\frac{16}{9}\xi$~eV, später mit
Korrektur zweiter Ordnung (P11):
\begin{equation}
k_BT_{\mathrm{CMB}}=\tfrac{16}{9}\,\xi\,\bigl(1-\tfrac{275}{4}\xi\bigr)\ \mathrm{eV}.
\end{equation}
In SI ergibt die erste Ordnung $3{,}7978\times10^{-23}$~J, also $2{,}7507$~K und $+0{,}93\,\%$; die
zweite Ordnung $2{,}72549$~K und $+0{,}0002\,\%$ \KM. Die Übereinstimmung ist echt, ihre Bedeutung
nicht: Drei Beobachtungen zeigen, dass sie an der Einheit hängt.

\emph{Erstens} stimmt die Relation nur in eV. Liest man $\frac{16}{9}\xi$ in meV, in Kelvin oder
in Einheiten von $m_e$, liegt sie um Faktoren $10^3$, $10^4$ oder $5\times10^5$ daneben \KM. Als
Verhältnis geschrieben lautet sie $T_{\mathrm{CMB}}/m_e=\frac{16}{9}\xi/510\,999$; die Zahl
$510\,999$ ist $m_e$ in eV, also der Anker selbst, und kein Ausdruck in $\xi$.

\emph{Zweitens} ändert die Umdefinition von $\alpha$ daran nichts. Mit $\alpha=1$ wird die
Ladungseinheit um $1/\sqrt\alpha=11{,}706$ größer und das Potential um denselben Faktor kleiner;
Energien bleiben gleich \BM. In $T_{\mathrm{CMB}}$ steckt keine Ladung, also kann $\alpha$ die
Relation weder erzeugen noch beseitigen.

\emph{Drittens} ist das eV gesetzt, nicht gemessen \QM. Seine Größe geht auf das praktische Volt
von 1881 zurück ($10^8$ elektromagnetische CGS-Einheiten, gewählt nach der Spannung der
Daniell-Zelle), damit auf Zentimeter, Gramm und Sekunde. Jede spätere Stufe -- internationales
Volt 1908, absolute Einheiten 1948, Josephson-Konstante 1990 -- hat diese Größe beibehalten. Seit
dem 20.~Mai 2019 gilt mit der festgesetzten Elementarladung exakt
$1\ \mathrm{eV}=1{,}602176634\times10^{-19}$~J (SI-Broschüre, 9.~Aufl.). Die CMB wurde 1965
entdeckt; in die Größe des eV kann sie nicht eingegangen sein.

Dazu kommt der Koeffizient. Exakt verlangt wäre $68{,}766$; $275/4=68{,}75$ ist der nächste Bruch
mit Nenner~4. Brüche $n/4$ liegen in $T$ um $3{,}3\times10^{-5}$ auseinander, ein passender findet
sich also immer auf $\pm1{,}7\times10^{-5}$. Die zweite Ordnung liegt $0{,}01\sigma$ neben dem
Mittelwert. Über die Messgenauigkeit hinaus belegt ein Treffer nichts; eine zutreffende Formel läge
typischerweise um $1\sigma$ daneben, so nahe trifft sie nur mit rund $1\,\%$ Wahrscheinlichkeit.
Das spricht dafür, dass der Koeffizient am Mittelwert gewählt wurde.
Die Einstufung von P11 in R137~(i) bleibt: einheitenabhängige Übereinstimmung, Koeffizient nicht
hergeleitet \SM.

Für die Buchführung heißt das: Eine Formel der Form \glqq Größe $=$ Zahl mal eV\grqq{} benutzt die
Einheit eV als Anker. Das ist erlaubt, solange es der einzige ist. Die Massenleiter benutzt aber
$m_e$; beide zusammen verlangen, dass $m_e/\mathrm{eV}$ aus $\xi$ folgt. Die einzige Nähe ist
$\mathrm{eV}/m_e=1{,}271\,\xi^{3/2}$, $0{,}17\,\%$ neben $(4/\pi)\,\xi^{3/2}$ \KM; das würde die
Größe des Volt aus $\xi$ ableiten und ist nicht tragfähig \XM.

% ============================================================
\section{Andere Bezugsenergien}
% ============================================================

Die Grundform legt selbst eine Energieeinheit fest: Aus $\xi E_0^2=1$ folgt
$u=\sqrt{\xi}\,E_0$, mit $E_0=\sqrt{m_em_\mu}$ also $u=84{,}85$~keV. In dieser Einheit ist
$T_{\mathrm{CMB}}=2{,}768\times10^{-9}\,u=0{,}0876\cdot\frac{16}{9}\xi^2\,u$ \KM, ohne glatten
Vorfaktor. Auf die Energieskala des Atoms bezogen gilt
$T_{\mathrm{CMB}}/(\alpha^2m_e)=0{,}0647\,\xi$ \KM, ebenfalls ohne glatten Vorfaktor. Beide
Bezüge liefern keine Relation.

% ============================================================
\section{Die Grundform: $T_{\mathrm{CMB}}/m_e$}
% ============================================================

In natürlichen Einheiten hängt $T_{\mathrm{CMB}}$ an der Hubble-Rate, der Planck-Energie und dem
Anteil der Photonen an der kritischen Dichte. Mit $\rho_\gamma=\frac{\pi^2}{15}T^4$ und
$\rho_{\mathrm{krit}}=3H_0^2E_P^2/(8\pi)$ gilt
\begin{equation}
\frac{T_{\mathrm{CMB}}}{m_e}=\Bigl(\frac{45\,\Omega_\gamma}{8\pi^3}\Bigr)^{1/4}
\Bigl(\frac{H_0E_P}{m_e^2}\Bigr)^{1/2}.
\end{equation}
Mit den Formen des Korpus $H_0=\frac{\pi}{2}\xi^{10}m_e$ (Dok.~165, 279) und
$m_e/E_P=\frac{4/3}{5\pi}\xi^{11/2}$ (Dok.~383) wird $H_0E_P/m_e^2=\frac{15\pi^2}{8}\xi^{9/2}$
\BM. Ist $\Omega_\gamma$ proportional zu $\xi$, folgt $T_{\mathrm{CMB}}/m_e\propto\xi^{5/2}$.
Gemessen ist
\begin{equation}
\frac{T_{\mathrm{CMB}}}{m_e}=K\,\xi^{5/2},\qquad K=2{,}23897 \quad\KM.
\end{equation}
Das erklärt den Exponenten $2{,}41$, den man ohne Vorfaktor abliest. Die Frage ist damit auf eine
Zahl zurückgeführt, gleichbedeutend mit $\Omega_\gamma/\xi$.

Einfache Kandidaten liegen außerhalb der Messgenauigkeit: $\sqrt5$ ($-0{,}13\,\%$), $9/4$
($+0{,}49\,\%$), $64/(9\pi)$ ($+1{,}10\,\%$) \XM. Innerhalb liegt
\begin{equation}
\boxed{\frac{T_{\mathrm{CMB}}}{m_e}=(8\pi)^{1/4}\,\xi^{5/2}}
\qquad K=2{,}23903,\ +0{,}0025\,\%,\ 0{,}12\sigma \quad\KM.
\end{equation}
$8\pi$ ist die Kopplung in $8\pi G$, was den Kandidaten motiviert, aber nicht begründet \SM.

Zur Zufallsprüfung wurden alle $3370$ Ausdrücke $(p/q\cdot\pi^k)^{1/n}$ mit $p,q\le30$,
$|k|\le3$, $n\le4$ im Bereich $1{,}5$ bis $3{,}5$ durchsucht. Drei treffen innerhalb von FIRAS,
zufällig erwartet sind $1{,}6$ \KM. Beschränkt auf $p,q\le10$ (383 Ausdrücke) bleibt nur
$(8\pi)^{1/4}$, bei $0{,}2$ zufällig erwarteten Treffern. Der Kandidat ist damit der einfachste,
aber nicht der einzige; ein Zufall ist nicht ausgeschlossen.

% ============================================================
\section{Folgerungen des Kandidaten}
% ============================================================

\emph{Beziehung zu $H_0$.} Gleichwertig gilt $T_{\mathrm{CMB}}^4=8\pi\,\xi^{10}m_e^4$, mit
demselben Exponenten~10 wie in $H_0$. Zusammen mit $H_0=\frac{\pi}{2}\xi^{10}m_e$ verschwindet
$\xi$:
\begin{equation}
\boxed{T_{\mathrm{CMB}}^4=16\,H_0\,m_e^3}\qquad\BM.
\end{equation}
Aus den gemessenen $T_{\mathrm{CMB}}$ und $m_e$ folgt $H_0=66{,}81\pm0{,}06$~km/s/Mpc \KM. Das
liegt $1{,}2\sigma$ neben Planck ($67{,}4\pm0{,}5$) und weit neben SH0ES. Die Unsicherheit kommt
allein aus FIRAS und ist rund achtmal kleiner als die der direkten $H_0$-Messungen.

\emph{Photonendichte.} Mit den T0-Werten von $H_0$ und $E_P$ wird
$\Omega_\gamma=\frac{4096}{10125}\,\xi=0{,}4045\,\xi$ \BM; mit dem gemessenen $E_P$ statt der
T0-Form $0{,}4154\,\xi$.

\emph{$L_\xi$ ohne CMB.} Setzt man den Kandidaten in die Definition von $L_\xi$ ein, fällt
$T_{\mathrm{CMB}}$ heraus:
\begin{equation}
L_\xi=\Bigl(\frac{15}{8\pi^3}\Bigr)^{1/4}\xi^{-9/4}\,\bar\lambda_e=100{,}24\,\mu\mathrm{m}\qquad\KM.
\end{equation}
$L_\xi$ ist damit eine Länge aus $\xi$ und der Compton-Wellenlänge des Elektrons. Das Verhältnis
$\pi^2/(720\xi)=102{,}8$ zwischen Casimir- und CMB-Dichte bei $L_\xi$ bleibt eine Identität; die
Zirkularität verschiebt sich aber von $L_\xi$ auf den Vorfaktor $(8\pi)^{1/4}$.

% ============================================================
\section{Die $H_0$-Formen im Vergleich}
% ============================================================

Im Korpus stehen zwei Formen: $E_H=E_0\,\xi^{41/4}$ (Dok.~026, 064; Exponent am Messwert
kalibriert, P35) und $H_0=\frac{\pi}{2}\xi^{10}m_e$ (Dok.~165, 279). Sie unterscheiden sich um
$1\,\%$, denn $E_0\,\xi^{1/4}=0{,}9904\cdot\frac{\pi}{2}m_e$ \KM. In der $H_0$-Messung fällt das nicht
auf; $T_{\mathrm{CMB}}$ ist aber mehr als dreißigmal genauer gemessen und trennt die Formen über
$T^4=16\,H_0m_e^3$:

\begin{center}
\small
\begin{tabular}{@{}p{5.2cm}p{2.0cm}p{2.0cm}p{3.6cm}@{}}
\toprule
$H_0$-Form & $H_0$ (km/s/Mpc) & $T_{\mathrm{CMB}}$ (K) & Abweichung \\
\midrule
$\frac{\pi}{2}\xi^{10}m_e$ & $66{,}821$ & $2{,}72555$ & $+0{,}003\,\%$, $0{,}1\sigma$ \KM \\
$E_0\,\xi^{41/4}$, $E_0=7{,}398$~MeV & $66{,}179$ & $2{,}71898$ & $-0{,}24\,\%$, $11\sigma$ \XM \\
$E_0\,\xi^{41/4}$, $E_0=\sqrt{m_em_\mu}$ & $65{,}731$ & $2{,}71436$ & $-0{,}41\,\%$, $20\sigma$ \XM \\
aus $T_{\mathrm{CMB}}$ rückwärts & $66{,}81\pm0{,}06$ & Messwert & -- \\
\bottomrule
\end{tabular}
\end{center}

Gilt der Kandidat, ist die Form mit $41/4$ durch die CMB ausgeschlossen, und der Übergang zum
Exponenten~10 in Dok.~165 ist bestätigt. Unabhängig ist diese Bestätigung nicht: Der Vorfaktor
$(8\pi)^{1/4}$ wurde am Messwert gesucht, und die Formel enthält $\xi^{10}$ bereits. Gezeigt ist,
dass dieselbe $\xi^{10}$-Struktur $H_0$ und $T_{\mathrm{CMB}}$ zugleich trägt, mit den einfachen
Vorfaktoren $\pi/2$ und $8\pi$.

% ============================================================
\section{Gegenüberstellung}
% ============================================================

\begin{center}
\small
\begin{tabular}{@{}p{4.4cm}p{2.6cm}p{2.0cm}p{4.2cm}@{}}
\toprule
Weg & Anker & Treffer & Einstufung \\
\midrule
$\frac{16}{9}\xi$~eV & eV & $+0{,}93\,\%$ & einheitenabhängig \SM \\
$\frac{16}{9}\xi(1-\frac{275}{4}\xi)$~eV & eV & $2\times10^{-6}$ & einheitenabhängig, angepasst \SM \\
$T/u$, $u=\sqrt\xi E_0$ & $m_e$ & kein Vorfaktor & keine Relation \\
$T/(\alpha^2m_e)$ & $m_e$ & kein Vorfaktor & keine Relation \\
$\sqrt5$, $9/4$, $64/(9\pi)$ & $m_e$ & $0{,}13$ bis $1{,}1\,\%$ & außerhalb FIRAS \XM \\
$(8\pi)^{1/4}\xi^{5/2}$ & $m_e$ & $0{,}12\sigma$ & Kandidat \SM \\
$\Lambda$CDM & $T_{\mathrm{CMB}}$ selbst & -- & Messwert als Eingabe \QM \\
\bottomrule
\end{tabular}
\end{center}

% ============================================================
\section{Was offen bleibt}
% ============================================================

Gelöst ist die Zirkularität von $L_\xi$ nur bedingt: Sie ist auf eine Zahl verschoben. Offen sind
\begin{itemize}
\item der Vorfaktor $(8\pi)^{1/4}$, gleichbedeutend mit $\Omega_\gamma=\frac{4096}{10125}\xi$:
motiviert durch $8\pi G$, nicht hergeleitet; ein Zufall ist bei $1{,}6$ erwarteten Treffern im
weiten Suchraum nicht ausgeschlossen \SM;
\item der Exponent~10 und der Vorfaktor $\pi/2$ von $H_0$ (P20); die $T$-Relation erbt sie, denn
$5/2=10/4$ \SM;
\item warum die Photonendichte heute ein fester Bruchteil von $\xi$ ist; in einem expandierenden
Modell wäre das ein Zeitpunktzufall, in einem statischen eine Aussage, die eine Herleitung braucht
\SM;
\item die eV-Relation P11: einheitenabhängig, Koeffizient $275/4$ angepasst; sie bleibt als
historische Übereinstimmung stehen, trägt aber nichts zur Grundform bei (R137, R70) \SM.
\end{itemize}
Prüfbar ist der Kandidat über $H_0$: Er verlangt $H_0=66{,}81\pm0{,}06$~km/s/Mpc. Eine
modellunabhängige $H_0$-Messung mit wenigen Zehntel Prozent Genauigkeit würde ihn bestätigen
oder ausschließen.

% ============================================================
\section{Fazit}
% ============================================================

Die eV-Relation trifft, weil das eV die Elektronenmasse als zweiten Anker mitbringt; in jeder
Einheit, die die FFGFT selbst festlegt, verschwindet die Übereinstimmung. In der Grundform hängt
$T_{\mathrm{CMB}}/m_e$ an einer einzigen Zahl, und der einfachste Kandidat,
$(8\pi)^{1/4}\xi^{5/2}$, trifft FIRAS auf $0{,}12\sigma$. Er verbindet $T_{\mathrm{CMB}}$ und $H_0$
zu $T^4=16\,H_0m_e^3$, legt $L_\xi$ ohne CMB fest und schließt die alte $H_0$-Form mit $41/4$ aus.
Hergeleitet ist er nicht; offen bleiben der Vorfaktor und, wie für $H_0$, der Exponent~10.

\vspace{1em}
\noindent\textit{Prüfskript:} \texttt{2/python/Dok388\_Skripte/pruef\_388\_cmb\_grundform.py}
--- 27/27.
